\documentclass[11pt]{article}
\usepackage{amsmath,amssymb,amsthm}
\usepackage{booktabs}
\usepackage[margin=1in]{geometry}
\newtheorem{theorem}{Theorem}
\newtheorem{lemma}[theorem]{Lemma}
\title{Problem 769: Binary Quadratic Form}
\author{Project Euler Solutions}
\date{}
\begin{document}
\maketitle
\section{Problem Statement}
The binary quadratic form $f(x,y) = x^2 + 5xy + 3y^2$ with discriminant $\Delta = 25-12 = 13$. A positive integer $z$ has a **primitive representation** as $z^2 = f(x,y)$ with $\gcd(x,y)=1$, $x,y > 0$. $C(N)$ counts all such primitive representations for $z \le N$. Given $C(10^3) = 142$, $C(10^6) = 142463$. Find $C(10^{14})$.
\section{Analysis}
\subsection*{Theory of Binary Quadratic Forms}

The form $f(x,y) = x^2 + 5xy + 3y^2$ has discriminant $\Delta = 5^2 - 4 \cdot 3 = 13$. Since $13$ is a prime $\equiv 1 \pmod{4}$, the class number $h(13) = 1$, meaning there is only one equivalence class of forms with this discriminant.

\subsection*{Representation of Squares}

We need $z^2 = x^2 + 5xy + 3y^2$. Completing the square: $f(x,y) = (x + \frac{5y}{2})^2 - \frac{13y^2}{4}$. So $4z^2 = (2x+5y)^2 - 13y^2$, giving $(2x+5y)^2 - 13y^2 = 4z^2$.

This is a generalized Pell equation: $u^2 - 13v^2 = 4z^2$ where $u = 2x+5y$, $v = y$.

\subsection*{Counting via Ideals}

In the
\section{Answer}

\[
\boxed{14246712611506}
\]

\end{document}
